\documentclass[10pt,a4paper]{amsart}
\usepackage{amsmath,amssymb,amsthm,graphicx}
% For dealing with references we use the comment environment
\usepackage{verbatim}
\newenvironment{reference}{\comment}{\endcomment}
%\newenvironment{reference}{}{}
\newenvironment{slogan}{\comment}{\endcomment}
\newenvironment{history}{\comment}{\endcomment}
% For commutative diagrams we use Xy-pic
\usepackage[all]{xy}

% We use 2cell for 2-commutative diagrams.
\xyoption{2cell}
\UseAllTwocells
\theoremstyle{plain}
\newtheorem{theorem}[subsection]{Theorem}
\newtheorem{proposition}[subsection]{Proposition}
\newtheorem{lemma}[subsection]{Lemma}

\theoremstyle{definition}
\newtheorem{definition}[subsection]{Definition}
\newtheorem{example}[subsection]{Example}
\newtheorem{exercise}[subsection]{Exercise}
\newtheorem{situation}[subsection]{Situation}

\theoremstyle{remark}
\newtheorem{remark}[subsection]{Remark}
\newtheorem{remarks}[subsection]{Remarks}

\numberwithin{equation}{subsection}
% Macros
%
\def\lim{\mathop{\mathrm{lim}}\nolimits}
\def\colim{\mathop{\mathrm{colim}}\nolimits}
\def\Spec{\mathop{\mathrm{Spec}}}
\def\Hom{\mathop{\mathrm{Hom}}\nolimits}
\def\Ext{\mathop{\mathrm{Ext}}\nolimits}
\def\SheafHom{\mathop{\mathcal{H}\!\mathit{om}}\nolimits}
\def\SheafExt{\mathop{\mathcal{E}\!\mathit{xt}}\nolimits}
\def\Sch{\mathit{Sch}}
\def\Mor{\mathop{\mathrm{Mor}}\nolimits}
\def\Ob{\mathop{\mathrm{Ob}}\nolimits}
\def\Sh{\mathop{\mathit{Sh}}\nolimits}
\def\NL{\mathop{N\!L}\nolimits}
\def\CH{\mathop{\mathrm{CH}}\nolimits}
\def\proetale{{pro\text{-}\acute{e}tale}}
\def\etale{{\acute{e}tale}}
\def\QCoh{\mathit{QCoh}}
\def\Ker{\mathop{\mathrm{Ker}}}
\def\Im{\mathop{\mathrm{Im}}}
\def\Coker{\mathop{\mathrm{Coker}}}
\def\Coim{\mathop{\mathrm{Coim}}}

% Boxtimes
%
\DeclareMathSymbol{\boxtimes}{\mathbin}{AMSa}{"02}

%
% Macros for moduli stacks/spaces
%
\def\QCohstack{\mathcal{QC}\!\mathit{oh}}
\def\Cohstack{\mathcal{C}\!\mathit{oh}}
\def\Spacesstack{\mathcal{S}\!\mathit{paces}}
\def\Quotfunctor{\mathrm{Quot}}
\def\Hilbfunctor{\mathrm{Hilb}}
\def\Curvesstack{\mathcal{C}\!\mathit{urves}}
\def\Polarizedstack{\mathcal{P}\!\mathit{olarized}}
\def\Complexesstack{\mathcal{C}\!\mathit{omplexes}}
% \Pic is the operator that assigns to X its picard group, usage \Pic(X)
% \Picardstack_{X/B} denotes the Picard stack of X over B
% \Picardfunctor_{X/B} denotes the Picard functor of X over B
\def\Pic{\mathop{\mathrm{Pic}}\nolimits}
\def\Picardstack{\mathcal{P}\!\mathit{ic}}
\def\Picardfunctor{\mathrm{Pic}}
\def\Deformationcategory{\mathcal{D}\!\mathit{ef}}
\usepackage{cleveref}
\crefname{originalreference}{original reference}{original references}
\Crefname{originalreference}{Original reference}{Original references}
\makeatletter
\newlabel{example-schemes}{{0DY7}{source}{}{sourcepointer-1}{}}
\newlabel{example-schemes@cref}{{[originalreference][0][]0DY7}{[1][0][]source}{}{}{}}
\makeatother
\crefname{originalsection}{original section}{original sections}
\Crefname{originalsection}{Original section}{Original sections}
\makeatletter
\newlabel{more-morphisms-section-thickenings}{{04EW}{source}{}{sourcepointer-2}{}}
\newlabel{more-morphisms-section-thickenings@cref}{{[originalsection][0][]04EW}{[1][0][]source}{}{}{}}
\makeatother
\crefname{originallemma}{original lemma}{original lemmas}
\Crefname{originallemma}{Original lemma}{Original lemmas}
\makeatletter
\newlabel{more-morphisms-lemma-deform}{{063Y}{source}{}{sourcepointer-3}{}}
\newlabel{more-morphisms-lemma-deform@cref}{{[originallemma][0][]063Y}{[1][0][]source}{}{}{}}
\makeatother
\crefname{originallemma}{original lemma}{original lemmas}
\Crefname{originallemma}{Original lemma}{Original lemmas}
\makeatletter
\newlabel{defos-lemma-deform-module}{{08LI}{source}{}{sourcepointer-4}{}}
\newlabel{defos-lemma-deform-module@cref}{{[originallemma][0][]08LI}{[1][0][]source}{}{}{}}
\makeatother
\crefname{originallemma}{original lemma}{original lemmas}
\Crefname{originallemma}{Original lemma}{Original lemmas}
\makeatletter
\newlabel{more-morphisms-lemma-first-order-thickening}{{05YV}{source}{}{sourcepointer-5}{}}
\newlabel{more-morphisms-lemma-first-order-thickening@cref}{{[originallemma][0][]05YV}{[1][0][]source}{}{}{}}
\makeatother
\title{Lemma schemes morphisms smooth to base in examples-defos.tex}
\author{The Stacks Project Authors}
\date{stacks-project:a04446e57ec1fbc252a871afcec7752fb2807b14; GNU Free Documentation License 1.2 (as published in the Stacks Project COPYING)}
\begin{document}
\maketitle
\noindent\textit{Editorial note.} All mathematical statement and proof text is reproduced verbatim from the author's chapter examples-defos.tex of the Stacks Project, fixed at commit a04446e57ec1fbc252a871afcec7752fb2807b14. The pinned preamble is supplied as its own file; the wrapper class, the xr-hyper wiring and the \textbackslash{}externaldocument list are omitted because the assembler supplies the class.\par\medskip
\noindent\textit{Editorial note.} Numbering inside this standalone excerpt is editorial, not the full Stacks book print numbering.\par\medskip
\noindent\textit{Editorial note.} The candidate is from the extractor's stacks-candidates line, not from the earlier judged initial\_corpus line; the H2 judgement and its reason below were made in this module by reading the author's own statement and complete proof in the fixed checkout. Original candidate key: stacks:examples-defos.tex:lemma-schemes-morphisms-smooth-to-base.\par\medskip
\noindent\textit{Editorial note.} Reference labels in the proof resolve to navigation pointers to the original author source; these pointers do not reproduce proofs of the prerequisites.\par\medskip
\noindent\textit{Allowed standard prerequisites (editorial).}\begin{itemize}
\item Prerequisites named by the author inside the counted span. The counted span is reproduced verbatim, including the author's own citations. Its prerequisites are the earlier results of the same fixed Stacks Project version that the author cites there and in the chapter's lead-in.
\end{itemize}\medskip
\ifdefined\hypertarget\hypertarget{sourcepointer-1}{}\fi
\noindent\textit{Source reference.} Original reference 0DY7: sources/examples-defos.tex, source lines 1208--1208. Navigation pointer used in the counted span; the span contains the pointer only and the cited result is not reproduced inside it. Registered as a non-load-bearing navigation stub in the established form for Stacks material. No bibliography citation is involved.\par\medskip
\ifdefined\hypertarget\hypertarget{sourcepointer-2}{}\fi
\noindent\textit{Source reference.} Original section 04EW: sources/more-morphisms.tex, source lines 30--30. Navigation pointer used in the counted span; the span contains the pointer only and the cited result is not reproduced inside it. Registered as a non-load-bearing navigation stub in the established form for Stacks material. No bibliography citation is involved.\par\medskip
\ifdefined\hypertarget\hypertarget{sourcepointer-3}{}\fi
\noindent\textit{Source reference.} Original lemma 063Y: sources/more-morphisms.tex, source lines 2182--2182. Navigation pointer used in the counted span; the span contains the pointer only and the cited result is not reproduced inside it. Registered as a non-load-bearing navigation stub in the established form for Stacks material. No bibliography citation is involved.\par\medskip
\ifdefined\hypertarget\hypertarget{sourcepointer-4}{}\fi
\noindent\textit{Source reference.} Original lemma 08LI: sources/defos.tex, source lines 1319--1319. Navigation pointer used in the counted span; the span contains the pointer only and the cited result is not reproduced inside it. Registered as a non-load-bearing navigation stub in the established form for Stacks material. No bibliography citation is involved.\par\medskip
\ifdefined\hypertarget\hypertarget{sourcepointer-5}{}\fi
\noindent\textit{Source reference.} Original lemma 05YV: sources/more-morphisms.tex, source lines 92--92. Navigation pointer used in the counted span; the span contains the pointer only and the cited result is not reproduced inside it. Registered as a non-load-bearing navigation stub in the established form for Stacks material. No bibliography citation is involved.\par\medskip
\begin{lemma}
\label{lemma-schemes-morphisms-smooth-to-base}
\begin{reference}
This is discussed in \cite[Section 5.3]{Ravi-Murphys-Law} and
\cite[Theorem 3.3]{Ran-deformations}.
\end{reference}
In Example \ref{example-schemes} let $f : X \to Y$ be a morphism of schemes
over $k$. If $f_*\mathcal{O}_X = \mathcal{O}_Y$ and $R^1f_*\mathcal{O}_X = 0$,
then the morphism of deformation categories
$$
\Deformationcategory_{X \to Y} \to \Deformationcategory_X
$$
is an equivalence.
\end{lemma}

\begin{proof}
We construct a quasi-inverse to the forgetful functor of the lemma.
Namely, suppose that $(A, U)$ is an object of $\Deformationcategory_X$.
The given map $X \to U$ is a finite order thickening and we can use
it to identify the underlying topological spaces of $U$ and $X$, see
More on Morphisms, Section \ref{more-morphisms-section-thickenings}.
Thus we may and do think of $\mathcal{O}_U$ as a sheaf of
$A$-algebras on $X$; moreover the fact that $U \to \Spec(A)$ is
flat, means that $\mathcal{O}_U$ is flat as a sheaf of $A$-modules.
In particular, we have a filtration
$$
0 = \mathfrak m_A^n\mathcal{O}_U \subset
\mathfrak m_A^{n - 1}\mathcal{O}_U \subset \ldots \subset
\mathfrak m_A^2\mathcal{O}_U \subset
\mathfrak m_A\mathcal{O}_U \subset \mathcal{O}_U
$$
with subquotients equal to
$\mathcal{O}_X \otimes_k \mathfrak m_A^i/\mathfrak m_A^{i + 1}$
by flatness, see More on Morphisms, Lemma \ref{more-morphisms-lemma-deform}
or the more general Deformation Theory, Lemma \ref{defos-lemma-deform-module}.
Set
$$
\mathcal{O}_V = f_*\mathcal{O}_U
$$
viewed as sheaf of $A$-algebras on $Y$. Since
$R^1f_*\mathcal{O}_X = 0$ we find by the description above that
$R^1f_*(\mathfrak m_A^i\mathcal{O}_U/\mathfrak m_A^{i + 1}\mathcal{O}_U) = 0$
for all $i$. This implies that the sequences
$$
0 \to
(f_*\mathcal{O}_X) \otimes_k \mathfrak m_A^i/\mathfrak m_A^{i + 1} \to
f_*(\mathcal{O}_U/\mathfrak m_A^{i + 1}\mathcal{O}_U) \to
f_*(\mathcal{O}_U/\mathfrak m_A^i\mathcal{O}_U) \to 0
$$
are exact for all $i$. Reading the references given above backwards
(and using induction) we find that $\mathcal{O}_V$ is a flat
sheaf of $A$-algebras with
$\mathcal{O}_V/\mathfrak m_A\mathcal{O}_V = \mathcal{O}_Y$.
Using More on Morphisms, Lemma
\ref{more-morphisms-lemma-first-order-thickening}
we find that $(Y, \mathcal{O}_V)$ is a scheme, call it $V$.
The equality $\mathcal{O}_V = f_*\mathcal{O}_U$ defines a
morphism of ringed spaces $U \to V$ which is easily seen to be
a morphism of schemes. This finishes the proof by the
flatness already established.
\end{proof}

\end{document}
